\chapter{Transforms, Interpolation and Algorithmic Differentiation}\label{ch:qm:transforms-interpolation-and-algorithmic-differentiation}

The risk run bumps 400 curve points one at a time and reprices the book 401 times every night. The same 400 sensitivities come out of one recording of the pricing code and one sweep backwards
through it, at the cost of a few prices. This chapter collects the numerical tools that pricing and risk code is made of: Fourier methods that price from a \omterm{def:qm:probability-at-speed:charfn}{characteristic function}, interpolation that
turns a handful of quotes into a curve, root finders that invert a pricing formula, and \omterm{def:qm:transforms-interpolation-and-algorithmic-differentiation:ad}{algorithmic differentiation}, which delivers exact derivatives of whatever the code computes.

\section{Fourier methods}

\begin{definition}[Discrete and fast Fourier transforms]\label{def:qm:transforms-interpolation-and-algorithmic-differentiation:dft}
The \emph{discrete Fourier transform}\index{discrete Fourier transform} of $x_0, \dots, x_{N-1}$ is $X_k = \sum_{j=0}^{N-1}x_je^{-2\pi\mathrm ijk/N}$, $k = 0, \dots, N - 1$. A \emph{fast Fourier
transform}\index{fast Fourier transform} computes it in $O(N\log N)$ operations instead of $O(N^2)$ by splitting it recursively into transforms of half the length (Cooley and Tukey, 1965).
\end{definition}

A Lévy model is specified by its \omterm{def:qm:jump-processes:triplet}{characteristic exponent} (chapter~6), not its density, so prices come from the \omterm{def:qm:probability-at-speed:charfn}{characteristic function} $\varphi(u) = \E[e^{\mathrm iu\ln(S_T/S_0)}]$. Carr and Madan (1999)
price a whole strip of strikes with one \omterm{def:qm:transforms-interpolation-and-algorithmic-differentiation:dft}{fast Fourier transform} of a damped call; the method below needs no damping and converges faster.

\begin{proposition}[Fourier inversion]\label{prop:qm:transforms-interpolation-and-algorithmic-differentiation:gp}
For a random variable $X$ with \omterm{def:qm:probability-at-speed:charfn}{characteristic function} $\varphi$ and a continuity point $k$ of its distribution, $\P(X > k) = \frac12 + \frac1\pi\int_0^\infty\mathrm{Re}\bigl[e^{-\mathrm iuk}\varphi(u)/(\mathrm iu)\bigr]\,du$
(Gil-Pelaez, 1951).
\end{proposition}

For a Merton jump diffusion ($\sigma = 15\%$, jumps at rate 0.5 a year with mean $-0.1$ and standard deviation 0.2 in log-price, $r = 3\%$, one year), the formula with the trapezoidal rule on $[0, 200]$ gives
$\P(S_T > 110) = 0.3503999$, against the Poisson mixture of lognormals to $3 \times 10^{-10}$.

\begin{definition}[COS method]\label{def:qm:transforms-interpolation-and-algorithmic-differentiation:cos}
The \emph{COS method}\index{COS method} (Fang and Oosterlee, 2008) truncates the density of $y = \ln(S_T/K)$ to an interval $[a, b]$, expands it in a cosine series whose coefficients are read off the
\omterm{def:qm:probability-at-speed:charfn}{characteristic function}, $A_k \approx \frac2{b-a}\mathrm{Re}[\varphi(\frac{k\pi}{b-a})e^{-\mathrm ik\pi a/(b-a)}]$, and integrates the payoff against each cosine in closed form, so that a price is
$e^{-rT}\sum_{k<N}{}'A_kV_k$ (the first term halved).
\end{definition}

\begin{proposition}[Convergence of the COS method]\label{prop:qm:transforms-interpolation-and-algorithmic-differentiation:cosconv}
For a density that is smooth on the truncation interval, the cosine coefficients decay faster than any power of $k$, and the error of the COS price decays exponentially in $N$ until the \omterm{def:qm:finite-difference-methods:consistent}{truncation error}
of $[a, b]$ dominates (Fang and Oosterlee, 2008; stated).
\end{proposition}

With $[a, b]$ set by ten standard deviations from the \omterm{def:qm:jump-processes:cumulant}{cumulants}, a Black--Scholes call is exact to $7 \times 10^{-15}$ with 64 terms, and the Merton call to $10^{-7}$ with 64 and $2 \times 10^{-12}$ with 128
(\cref{fig:qm:transforms-interpolation-and-algorithmic-differentiation:cos}). The firm prices calls from the put expansion and parity, which keeps the coefficients bounded. For a variance-gamma model,
where only simulation is available as a check, 256 terms give 9.5243 and 400\,000 simulated paths $9.528 \pm 0.022$.

\begin{omfigure}
\begin{tikzpicture}
\begin{semilogyaxis}[width=10cm,height=5cm,xlabel={number of cosine terms $N$},ylabel={price error},
  tick label style={font=\small},label style={font=\small},xmin=0,xmax=264,ymin=1e-15,ymax=10,grid=major,grid style={gray!25},
  legend cell align=left,legend style={font=\footnotesize,at={(0.97,0.97)},anchor=north east,fill=white,draw=none},table/col sep=comma]
\addplot[omDef,very thick,mark=*,mark size=1.5pt] table[x=N,y=bs]{figdata/methods/28-transforms-interpolation-and-algorithmic-differentiation/cos.csv};
\addplot[omThm,very thick,mark=square*,mark size=1.5pt] table[x=N,y=merton]{figdata/methods/28-transforms-interpolation-and-algorithmic-differentiation/cos.csv};
\legend{Black--Scholes ($\sigma = 20\%$),Merton jump diffusion}
\end{semilogyaxis}
\end{tikzpicture}

\omcaption{Error of the COS price of a one-year at-the-money call against the number of cosine terms (logarithmic error axis, linear $N$): exponential convergence down to rounding for Black--Scholes and
to the accuracy of the reference series for Merton. Data: the chapter's tutorial.}\label{fig:qm:transforms-interpolation-and-algorithmic-differentiation:cos}
\end{omfigure}

\section{Interpolation and splines}

\begin{definition}[Cubic, natural and monotone splines, Runge phenomenon]\label{def:qm:transforms-interpolation-and-algorithmic-differentiation:spline}
A \emph{cubic spline}\index{cubic spline} through $(x_i, y_i)$ is a piecewise cubic with continuous first and second derivatives; a \emph{natural cubic spline}\index{natural cubic spline} has zero second
derivative at both ends, and its second derivatives at the nodes solve a tridiagonal system. \emph{Monotone cubic interpolation}\index{monotone cubic interpolation} chooses the node slopes of a piecewise
cubic Hermite interpolant so that it is monotone wherever the data are. The \emph{Runge phenomenon}\index{Runge phenomenon} (Runge, 1901) is the growth, near the ends of the interval, of the error of
polynomial interpolation at equally spaced nodes as the degree increases.
\end{definition}

\begin{proposition}[Minimal curvature of the natural spline]\label{prop:qm:transforms-interpolation-and-algorithmic-differentiation:minimal}
Among all twice continuously differentiable functions through the data, the \omterm{def:qm:transforms-interpolation-and-algorithmic-differentiation:spline}{natural cubic spline} minimises $\int(f^{\prime\prime})^2$.
\end{proposition}

\begin{proof}
Let $s$ be the spline and $g = f - s$, which vanishes at the nodes. Integrating by parts on each interval, $\int s^{\prime\prime}g^{\prime\prime} = [s^{\prime\prime}g']_{x_0}^{x_n} - \int s^{\prime\prime\prime}g' = 0$, because $s^{\prime\prime}$ vanishes at the ends and $s^{\prime\prime\prime}$ is
constant on each interval, where $\int g' = 0$. Hence $\int(f^{\prime\prime})^2 = \int(s^{\prime\prime})^2 + \int(g^{\prime\prime})^2 \ge \int(s^{\prime\prime})^2$.
\end{proof}

\begin{proposition}[Fritsch--Carlson condition]\label{prop:qm:transforms-interpolation-and-algorithmic-differentiation:fc}
On an interval where the data increase with secant slope $\delta > 0$, the cubic Hermite interpolant with end slopes $m_0, m_1 \ge 0$ is monotone if $\alpha = m_0/\delta$ and $\beta = m_1/\delta$ satisfy
$\alpha^2 + \beta^2 \le 9$ (Fritsch and Carlson, 1980).
\end{proposition}

Minimal curvature is not what a rates desk wants. On the Treasury curve of 3 July 2023 (constant-maturity yields from one month to thirty years, treated as zero rates for the illustration), the natural
spline of the zero rates produces a forward curve that falls to 2.63\% at thirty years; the monotone interpolation of $rt$, whose derivative is the forward, keeps every forward between 3.32\% and 5.67\%;
linear interpolation of zero rates gives forwards that jump by up to 0.43 percentage points at the pillars (\cref{fig:qm:transforms-interpolation-and-algorithmic-differentiation:fwd}). Hagan and West (2006)
survey the methods used for curves. Global polynomials are worse still: through eleven equally spaced samples of $1/(1 + 25x^2)$ on $[-1, 1]$, the interpolating polynomial errs by 1.92 near the ends.

\begin{omfigure}
\begin{tikzpicture}
\begin{axis}[width=11cm,height=5.4cm,xlabel={maturity (years)},ylabel={instantaneous forward (\%)},
  tick label style={font=\small},label style={font=\small},xmin=0,xmax=30,ymin=2.4,ymax=6,grid=major,grid style={gray!25},
  legend columns=4,legend cell align=left,legend style={font=\footnotesize,at={(0.5,-0.3)},anchor=north,draw=none,fill=none,
  /tikz/every even column/.append style={column sep=6pt}},table/col sep=comma]
\addplot[black!50,thick] table[x=t,y=linear]{figdata/methods/28-transforms-interpolation-and-algorithmic-differentiation/forwards.csv};
\addplot[omThm,very thick] table[x=t,y=spline]{figdata/methods/28-transforms-interpolation-and-algorithmic-differentiation/forwards.csv};
\addplot[omDef,very thick] table[x=t,y=monotone]{figdata/methods/28-transforms-interpolation-and-algorithmic-differentiation/forwards.csv};
\addplot[only marks,mark=*,mark size=1.6pt,black] table[x=t,y=zero]{figdata/methods/28-transforms-interpolation-and-algorithmic-differentiation/pillars.csv};
\legend{linear zeros,natural spline,monotone ($rt$),zero-rate pillars}
\end{axis}
\end{tikzpicture}

\omcaption{Instantaneous forward curves implied by three interpolations of the Treasury curve of 3 July 2023, with the pillar yields (treated as zero rates). Data: FRED series DGS1MO to DGS30 (Board of
Governors of the Federal Reserve System, H.15).}\label{fig:qm:transforms-interpolation-and-algorithmic-differentiation:fwd}
\end{omfigure}

\section{Root finding}

\begin{definition}[Bisection, secant and Brent's methods]\label{def:qm:transforms-interpolation-and-algorithmic-differentiation:roots}
The \emph{bisection method}\index{bisection method} halves a bracket $[a, b]$ with $f(a)f(b) < 0$ at each step. The \emph{secant method}\index{secant method} replaces the derivative in \omterm{def:qm:numerical-optimisation-in-practice:newton}{Newton's method} by the
slope through the last two iterates. \emph{Brent's method}\index{Brent's method} (Brent, 1973) keeps a bracket and takes an inverse quadratic interpolation or secant step when it stays inside the bracket and
shrinks it fast enough, and a bisection step otherwise.
\end{definition}

\begin{proposition}[Convergence of the root finders]\label{prop:qm:transforms-interpolation-and-algorithmic-differentiation:rootconv}
Bisection gains one bit per step and always converges; the \omterm{def:qm:transforms-interpolation-and-algorithmic-differentiation:roots}{secant method} converges locally with order $(1 + \sqrt5)/2 \approx 1.618$ at a simple root; \omterm{def:qm:transforms-interpolation-and-algorithmic-differentiation:roots}{Brent's method} always converges, never much more slowly than bisection
(at most about the square of the bisection count), and superlinearly near a simple root (Brent, 1973).
\end{proposition}

Implied volatility is the standard test. For a one-year call ($S = 100$, $r = 3\%$, true volatility 25\%), bisection on $[0.001, 5]$ needs 36 steps at every strike. At the money, Newton from 20\% needs 3
iterations, the \omterm{def:qm:transforms-interpolation-and-algorithmic-differentiation:roots}{secant method} 4, Brent 5. At a strike of 400 the price is $2.5 \times 10^{-7}$ and the vega at 20\% is $8 \times 10^{-9}$: Newton's first step goes to a volatility of 3\,062\%, where the vega
underflows, and the \omterm{def:qm:transforms-interpolation-and-algorithmic-differentiation:roots}{secant method} also fails, while Brent converges in 18. The firm's rule is a bracketing method with a fast interior step, and Newton only from a guess known to be close.

\section{Forward and adjoint algorithmic differentiation}

\begin{definition}[Algorithmic differentiation, forward and reverse modes, dual numbers]\label{def:qm:transforms-interpolation-and-algorithmic-differentiation:ad}
\emph{Algorithmic differentiation}\index{algorithmic differentiation} computes derivatives of a function given as a program by applying the chain rule to its elementary operations, exactly up to
rounding. \emph{Forward mode}\index{forward mode} propagates, with each value, its derivative along one input direction; \emph{dual numbers}\index{dual number} $a + b\epsilon$ with $\epsilon^2 = 0$ implement it,
since $f(a + b\epsilon) = f(a) + f'(a)b\epsilon$ (Wengert, 1964). \emph{Reverse mode}\index{reverse mode} records the operations and then propagates adjoints $\bar v = \partial y/\partial v$ from the output back to
every input (Linnainmaa, 1976, from his 1970 master's thesis on rounding errors; Griewank, 2012, recounts several independent discoveries); in finance it is called \emph{adjoint mode}\index{adjoint mode}, or AAD.
\end{definition}

\begin{definition}[Tape, checkpointing]\label{def:qm:transforms-interpolation-and-algorithmic-differentiation:tape}
A \emph{tape}\index{tape} is the record, made during the forward evaluation, of each elementary operation's inputs and local partial derivatives. \emph{Checkpointing}\index{checkpointing} stores only some
intermediate states during the forward pass and recomputes the operations between them during the reverse pass, trading computation for memory.
\end{definition}

\begin{proposition}[Cheap gradient principle]\label{prop:qm:transforms-interpolation-and-algorithmic-differentiation:cheap}
For a scalar function of $n$ inputs evaluated with $E$ elementary operations, each with at most two inputs, the \omterm{def:qm:transforms-interpolation-and-algorithmic-differentiation:ad}{reverse mode} computes the whole gradient with a number of operations bounded by a small
constant times $E$, independent of $n$; \omterm{def:qm:transforms-interpolation-and-algorithmic-differentiation:ad}{forward mode} and finite differences need $n$ passes (Griewank and Walther, 2008).
\end{proposition}

\begin{proof}
The recording evaluates each operation once and computes at most two local partials; the reverse sweep performs one multiply-add per recorded (operation, input) pair. With $P \le 2E$ pairs, the total
is $E + 2P \le 5E$ in this count, whatever $n$. \omterm{def:qm:transforms-interpolation-and-algorithmic-differentiation:ad}{Forward mode} propagates one tangent per pass, so the full gradient needs $n$ passes.
\end{proof}

The firm's toy book has 400 zero-rate pillars on thirty years, 120 quarterly cash flows and 119 caplets priced by Black's formula on the three-month forwards: 2\,864 operations for one value. The reverse
sweep costs, in the count above, 3.58 evaluations; bumping costs 401 and \omterm{def:qm:transforms-interpolation-and-algorithmic-differentiation:ad}{forward mode} 1\,433 (\cref{fig:qm:transforms-interpolation-and-algorithmic-differentiation:cost}). The adjoint gradient agrees
with \omterm{def:qm:transforms-interpolation-and-algorithmic-differentiation:ad}{forward mode} to $2 \times 10^{-13}$ and with central differences to $9 \times 10^{-6}$ on sensitivities as large as 517 (the bumps' own error); 213 of the 400 are nonzero, since only the pillars
next to a cash-flow date matter under linear interpolation. The C++20 twin, timed in its acceptance test, returns the gradient of a similar 400-input function in less than twenty times one plain
evaluation (between five and seven in the runs made for this chapter, \omterm{def:qm:transforms-interpolation-and-algorithmic-differentiation:tape}{tape} allocation included).

\begin{omfigure}
\begin{tikzpicture}
\begin{loglogaxis}[width=10cm,height=5cm,xlabel={number of inputs $n$},ylabel={cost of the gradient (evaluations)},
  tick label style={font=\small},label style={font=\small},xmin=20,xmax=500,ymin=1,ymax=3000,grid=major,grid style={gray!25},
  xtick={25,50,100,200,400},xticklabels={25,50,100,200,400},
  legend columns=3,legend cell align=left,legend style={font=\footnotesize,at={(0.5,-0.3)},anchor=north,draw=none,fill=none,
  /tikz/every even column/.append style={column sep=8pt}},table/col sep=comma]
\addplot[black,very thick,mark=square*,mark size=1.5pt] table[x=n,y=bump]{figdata/methods/28-transforms-interpolation-and-algorithmic-differentiation/adcost.csv};
\addplot[omThm,very thick,mark=*,mark size=1.5pt] table[x=n,y=forward]{figdata/methods/28-transforms-interpolation-and-algorithmic-differentiation/adcost.csv};
\addplot[omDef,very thick,mark=triangle*,mark size=1.8pt] table[x=n,y=reverse]{figdata/methods/28-transforms-interpolation-and-algorithmic-differentiation/adcost.csv};
\legend{bumping ($n + 1$),forward mode,reverse mode}
\end{loglogaxis}
\end{tikzpicture}

\omcaption{Cost of the gradient of the toy rates book, in evaluations, by operation count, against the number of zero-rate pillars: bumping and \omterm{def:qm:transforms-interpolation-and-algorithmic-differentiation:ad}{forward mode} grow with $n$, the \omterm{def:qm:transforms-interpolation-and-algorithmic-differentiation:ad}{reverse mode} does not.
Data: the chapter's tutorial.}\label{fig:qm:transforms-interpolation-and-algorithmic-differentiation:cost}
\end{omfigure}

The price of the \omterm{def:qm:transforms-interpolation-and-algorithmic-differentiation:ad}{reverse mode} is memory: the \omterm{def:qm:transforms-interpolation-and-algorithmic-differentiation:tape}{tape} holds every operation. A time-stepping loop of 10\,000 steps records 60\,003 nodes; with a checkpoint every 100 steps, the reverse pass holds at most 201
states (100 checkpoints and one segment) and recomputes the loop once, and the gradient is the same to $10^{-14}$. Two practical cautions: a payoff with a kink or a jump has a derivative that is zero or
undefined almost everywhere, so AAD of a Monte Carlo digital returns zero and needs smoothing or a likelihood-ratio term; and code that branches on values records only the branch taken. Giles and
Glasserman (2006) brought the adjoint method to Monte Carlo Greeks in finance.

\section{Tutorial: one sweep instead of 401}\label{tut:qm:transforms-interpolation-and-algorithmic-differentiation:aad}

\begin{tutorial}
\textbf{Goal.} Price by COS from a Lévy exponent, build three forward curves, invert Black--Scholes with four root finders, and differentiate a 400-input book three ways. \textbf{End state:} the four figures
and the gradient agreement.
\begin{tutsteps}
\item \textbf{The \omterm{def:qm:transforms-interpolation-and-algorithmic-differentiation:tape}{tape}.} Nodes, the reverse sweep and one overloaded operation in the C++20 header.
\omcode{code/firm/aad/cpp/firm_aad.hpp}{13}{52}{A tape, its reverse sweep and a recorded multiplication (C++20).}
\item \textbf{\omterm{def:qm:transforms-interpolation-and-algorithmic-differentiation:tape}{Checkpointing}.} The Python reference recomputes each segment on a fresh \omterm{def:qm:transforms-interpolation-and-algorithmic-differentiation:tape}{tape}.
\omcode{code/firm/aad/firm_aad.py}{201}{229}{Reverse mode through a long loop with checkpoints (Python).}
\item \textbf{Run} \texttt{cos\_convergence()}, \texttt{vg\_check()}, \texttt{gil\_pelaez\_digital()}, \texttt{forward\_curves()}, \texttt{implied\_vol\_iterations()}, \texttt{gradients()} and
\texttt{checkpoint\_demo()} in \texttt{qm\_transforms.py}; then \texttt{fig\_transforms.py}.
\end{tutsteps}
\textbf{What to change next.} Replace linear by monotone interpolation inside the book and watch the number of nonzero sensitivities; differentiate the COS price with respect to the Merton parameters by
running \texttt{cos\_call} on \omterm{def:qm:transforms-interpolation-and-algorithmic-differentiation:ad}{dual numbers}; add a second-order (forward over reverse) gamma.
\end{tutorial}

\section{Build: algorithmic differentiation}\label{bld:qm:transforms-interpolation-and-algorithmic-differentiation:aad}

\begin{build}
\bfield{Purpose} Exact sensitivities of any pricing code at the cost of a few evaluations, for the pricing library and the risk engine of later books.
\bfield{Interface} Python reference: \texttt{Tape}, \texttt{Var}, \texttt{Dual}, \texttt{exp}, \texttt{log}, \texttt{sqrt}, \texttt{ncdf}, \texttt{maximum}, \texttt{gradient}, \texttt{forward\_gradient},
\texttt{bump\_gradient}, \texttt{checkpointed\_gradient}. C++20: \texttt{firm::aad::Var} on a thread-local \texttt{tape}, \texttt{Tape::reverse}, \texttt{Dual}. Rust: \texttt{Var}, \texttt{reverse},
\texttt{clear}, \texttt{Dual}, \texttt{ncdf}.
\bfield{Rules} Every AAD gradient ships with a bump check on a sample of inputs; the \omterm{def:qm:transforms-interpolation-and-algorithmic-differentiation:tape}{tape} is cleared per valuation and per thread; kinks and jumps are smoothed before differentiation; long loops are
checkpointed.
\bfield{Acceptance tests} \texttt{code/firm/aad/\{tests,cpp,rust\}}: elementary rules against hand derivatives; operation counts; the 400-input test function's value and three gradient components
equal in the three languages; forward against reverse to $10^{-11}$; central and one-sided bumps; \omterm{def:qm:transforms-interpolation-and-algorithmic-differentiation:tape}{checkpointing} equal to the full \omterm{def:qm:transforms-interpolation-and-algorithmic-differentiation:tape}{tape}; the C++ time ratio below twenty.
\bfield{Stretch} Vector \omterm{def:qm:transforms-interpolation-and-algorithmic-differentiation:ad}{forward mode}; second derivatives; expression templates to shrink the \omterm{def:qm:transforms-interpolation-and-algorithmic-differentiation:tape}{tape}; adjoints of a linear solve (chapter~25) without taping its inner loops.
\end{build}

\begin{omsources}
\item J.~W.~Cooley and J.~W.~Tukey, \emph{Mathematics of Computation} 19, 1965; P.~Carr and D.~Madan, ``Option valuation using the fast Fourier transform'', \emph{Journal of Computational Finance} 2(4), 1999.
\item F.~Fang and C.~W.~Oosterlee, ``A novel pricing method for European options based on Fourier-cosine series expansions'', \emph{SIAM Journal on Scientific Computing} 31, 2008; J.~Gil-Pelaez, \emph{Biometrika} 38, 1951.
\item F.~N.~Fritsch and R.~E.~Carlson, ``Monotone piecewise cubic interpolation'', \emph{SIAM Journal on Numerical Analysis} 17, 1980; P.~S.~Hagan and G.~West, ``Interpolation methods for curve construction'',
\emph{Applied Mathematical Finance} 13, 2006; C.~Runge, \emph{Zeitschrift für Mathematik und Physik} 46, 1901.
\item R.~P.~Brent, \emph{Algorithms for Minimization without Derivatives}, Prentice-Hall, 1973.
\item R.~E.~Wengert, \emph{Communications of the ACM} 7, 1964; S.~Linnainmaa, \emph{BIT} 16, 1976; A.~Griewank, ``Who invented the reverse mode of differentiation?'', \emph{Documenta Mathematica}, Extra Volume ISMP, 2012; A.~Griewank and A.~Walther, \emph{Evaluating Derivatives}, 2nd ed., SIAM, 2008; M.~Giles and P.~Glasserman,
``Smoking adjoints: fast Monte Carlo Greeks'', \emph{Risk}, January 2006.
\item Treasury yields: Board of Governors of the Federal Reserve System, H.15, via FRED (series DGS1MO to DGS30), accessed 24 September 2026.
\end{omsources}

\section{Exercises}

\begin{exercise}[$\star$]\label{exo:qm:transforms-interpolation-and-algorithmic-differentiation:1}
How many complex multiplications does a direct \omterm{def:qm:transforms-interpolation-and-algorithmic-differentiation:dft}{discrete Fourier transform} of length 4\,096 need, and roughly how many does a radix-2 \omterm{def:qm:transforms-interpolation-and-algorithmic-differentiation:dft}{fast Fourier transform} need?
\end{exercise}

\begin{exercise}[$\star$]\label{exo:qm:transforms-interpolation-and-algorithmic-differentiation:2}
Evaluate $f(x) = x^2e^x$ on the \omterm{def:qm:transforms-interpolation-and-algorithmic-differentiation:ad}{dual number} $1 + \epsilon$ and read off $f'(1)$.
\end{exercise}

\begin{exercise}[$\star$]\label{exo:qm:transforms-interpolation-and-algorithmic-differentiation:3}
How many bisection steps reduce a bracket of width 5 below $10^{-10}$?
\end{exercise}

\begin{exercise}[$\star\star$]\label{exo:qm:transforms-interpolation-and-algorithmic-differentiation:4}
Write the \omterm{def:qm:transforms-interpolation-and-algorithmic-differentiation:tape}{tape} of $y = \sin(x_1x_2) + x_1$ and run the reverse sweep by hand at $(x_1, x_2) = (1, 2)$.
\end{exercise}

\begin{exercise}[$\star\star$]\label{exo:qm:transforms-interpolation-and-algorithmic-differentiation:5}
Show that the instantaneous forward is $f(t) = r(t) + tr'(t)$, and explain why a spline of $r$ with a large negative slope at the long end produces low forwards there.
\end{exercise}

\begin{exercise}[$\star\star$]\label{exo:qm:transforms-interpolation-and-algorithmic-differentiation:6}
Why does a pathwise AAD delta of a Monte Carlo digital option return zero, and name two remedies.
\end{exercise}

\begin{exercise}[$\star\star\star$]\label{exo:qm:transforms-interpolation-and-algorithmic-differentiation:7}
\emph{Coding.} Record the implied-volatility iteration counts of the four root finders at strikes 50, 70, 100, 150, 250 and 400, and explain the failures.
\end{exercise}

\begin{exercise}[$\star\star\star$]\label{exo:qm:transforms-interpolation-and-algorithmic-differentiation:8}
\emph{Find the flaw.} ``Our AAD gradient differs from the bumped one by $10^{-5}$ on a sensitivity of 500, so the AAD library has a bug.''
\end{exercise}

\section{Problem: One Sweep Instead of 401}

\begin{problem}[{Weekend problem --- 400 sensitivities for the price of a few}]\label{pb:qm:transforms-interpolation-and-algorithmic-differentiation:1}
A toy rates book on a 400-pillar zero curve holds 120 quarterly cash flows and 119 caplets. The nightly risk run bumps each pillar.

\medskip\noindent\textbf{Part I --- Prices from transforms.}
\begin{enumerate}
\item How does the COS error behave with the number of terms for Black--Scholes and Merton?
\item What does the Gil-Pelaez formula give for the Merton digital, and against what?
\item How is the variance-gamma COS price checked?
\item Why price calls from the put expansion?
\item When is the \omterm{def:qm:transforms-interpolation-and-algorithmic-differentiation:dft}{fast Fourier transform} the better tool?
\end{enumerate}

\medskip\noindent\textbf{Part II --- Curves and roots.}
\begin{enumerate}[resume]
\item What happens to the forwards under the natural spline of zero rates, and where?
\item What does monotone interpolation of $rt$ guarantee, and what range of forwards does it give?
\item What is wrong with linear interpolation of zero rates?
\item How do the four root finders fare at the money and at a strike of 400?
\item What is the Runge error on eleven equispaced nodes?
\end{enumerate}

\medskip\noindent\textbf{Part III --- The gradient.}
\begin{enumerate}[resume]
\item How many operations does one valuation perform, and how many (operation, input) pairs?
\item What do bumping, \omterm{def:qm:transforms-interpolation-and-algorithmic-differentiation:ad}{forward mode} and \omterm{def:qm:transforms-interpolation-and-algorithmic-differentiation:ad}{reverse mode} cost, in evaluations?
\item How well do the three gradients agree, and why does the bump differ?
\item Why are only 213 of 400 sensitivities nonzero?
\item What does \omterm{def:qm:transforms-interpolation-and-algorithmic-differentiation:tape}{checkpointing} buy on a 10\,000-step loop?
\end{enumerate}

\medskip\noindent\textbf{Part IV --- Judgement.}
\begin{enumerate}[resume]
\item Which sensitivities should still be bumped?
\item What must the risk engine check every night?
\item What does the C++ timing add to the operation count?
\item State the \emph{named result}: the cost ratio of the adjoint gradient to one evaluation for the 400-input book, against 401 for bumping, and the agreement between the two.
\item In one sentence: why is the \omterm{def:qm:transforms-interpolation-and-algorithmic-differentiation:ad}{adjoint mode} the natural mode for risk?
\end{enumerate}
\end{problem}

\section{Interview questions}

\begin{interviewq}[$\star$ \iqroles{researcher, developer}]\label{iq:qm:transforms-interpolation-and-algorithmic-differentiation:1}
Explain forward and \omterm{def:qm:transforms-interpolation-and-algorithmic-differentiation:ad}{reverse mode} differentiation, and when each is cheaper.
\end{interviewq}

\begin{interviewq}[$\star\star$ \iqroles{developer}]\label{iq:qm:transforms-interpolation-and-algorithmic-differentiation:2}
How would you implement \omterm{def:qm:transforms-interpolation-and-algorithmic-differentiation:ad}{reverse mode} in C++, and what limits its use in production?
\end{interviewq}

\begin{interviewq}[$\star\star$ \iqroles{researcher}]\label{iq:qm:transforms-interpolation-and-algorithmic-differentiation:3}
How do you compute an implied volatility robustly for any strike?
\end{interviewq}

\begin{interviewq}[$\star\star$ \iqroles{researcher}]\label{iq:qm:transforms-interpolation-and-algorithmic-differentiation:4}
Which interpolation would you use for a discount curve, and why not a \omterm{def:qm:transforms-interpolation-and-algorithmic-differentiation:spline}{natural cubic spline} of zero rates?
\end{interviewq}

\begin{interviewq}[$\star\star$ \iqroles{researcher}]\label{iq:qm:transforms-interpolation-and-algorithmic-differentiation:5}
How would you price a European option under a model known only through its \omterm{def:qm:probability-at-speed:charfn}{characteristic function}?
\end{interviewq}

\begin{interviewq}[$\star\star\star$ \iqroles{developer, risk}]\label{iq:qm:transforms-interpolation-and-algorithmic-differentiation:6}
Your AAD Monte Carlo gives a zero delta for a barrier option. What happened and what do you do?
\end{interviewq}
